\begin{proof}
Starting at the top of the tower and proceeding downwards, apply the
applicable one of
\Cref{prop:p062-open-polydisc-bundle-descent,prop:p062-projective-bundle-descent,prop:p062-smooth-proper-split-toric-bundle-descent,prop:p062-proper-smooth-curve-bundle-descent}.
The hypotheses in the statement retain all geometric and field assumptions
of the proposition used at each stage.  The successive unique descents give
a unique \(b\in C^0(P,\mathbf R)\) such that
\[
                         h=\Pi^*b,
                         \qquad \dprimeddprime[b]=0.
\]

Restrict the last current equation to the absolute analytic interior of
\(P\).  The endpoint-polycylinder theorem
\Cref{thm:p062-endpoint-polycylinder-v4-research} gives unique real constants
\(c,\lambda_1,\ldots,\lambda_a\) and the formula
\[
                   b=c+\sum_{i=1}^a\lambda_i\log|S_i|
                   \qquad\text{on all of }P.
\]
This is exactly \eqref{eq:p213-endpoint-mixed-tower-classification}; the
uniqueness of every descent and of the endpoint coefficients proves the
stated uniqueness.

On the standard tropical chart of an annular factor, \(\log|S_i|\) is the
pullback of an affine coordinate.  Hence \(b\) is ordinary CLD-smooth with
zero smooth Hessian.  Pullback of ordinary CLD forms is functorial and
commutes with both differentials by \Cref{prop:cld-form-pullback}, so
\(h=\Pi^*b\) has the same two properties.  If there is no annular factor,
the displayed formula reduces to \(h=\Pi^*c=c\).
\end{proof}
